\chapter*{Appendix III}

\addcontentsline{toc}{chapter}{Appendix III}

\def\shorttitle{Appendix III}

\section*{Sequence of modelling tasks]{Sequence of modelling tasks}

Modelling the fate of naturally attenuating contaminats in
groundwater consists of a series of sequential steps. Of course,
each site has its own characteristics with respect to the:
\begin{itemize}
    \item contaminant composition/mixture
    \item contamination history and available details of this history
    \item availability of hard and soft data
    \item hydrogeology
    \item hydrochemistry and mineralogy
\end{itemize}
and thus the modelling methodology will need to be adapted to
individual, site-specific requirements. However, the steps listed
in Table \ref{checklist} might serve as a general
guideline/checklist that provides a basis for a site-specific
schedule of a modelling study.

  \begin{table}
  \centering
%  \begin{tabular}{cr@{--}l} \hline
  \begin{tabular}{cp{5in}} \hline
    Step & Activity  \\ \hline
    1 & Initial assessment of the
    existing data/situation to be
modelled - initial identification of potential physical and
chemical interactions including initial `hand' calculations of
flow rates and initial `electron balance', where possible \\
    2 & Compilation of a first conceptual model for the
contaminated site. In particular, detailed depth profile data
(from multi-level sampling devices) can help for-mulate the
conceptual model for a site. Care is needed as to the source
conditions for the plume source - NAPL concentrations, source
dimensions. They will strongly influence plume characteristics and
the longevity of the contamination source \\
    3 & Formulation of a modelling strategy (dimensionality,
spatial and temporal resolution, proc-ess-detail of modelled
physical and chemical processes) and formulation of the questions
that the model(s) will need to answer. When deciding upon the
process detail of the chemical processes, i.e., the reaction
model, remember that: � Components of petrol products, including
benzene, toluene, ethylbenzene, naphthalene, etc do not degrade at
similar rates or under similar geochemical conditions - for
exam-ple, often toluene will degrade more readily than benzene;
and � Zero-, first-order or any non-specific third-party reaction
rate should only be used where groundwater and aquifer mineral
geochemical signatures are similar.  \\
    4 & Selection of appropriate modelling tools,
depending on available data, questions to be an-swered and
available time. \\
    5 & Setup of a flow model, including spatial
discretisation of the model domain into grid cells, definition of
boundary conditions, temporal discretisation of the simulation
period, e.g., to account for seasonal recharge dynamics, time
stepping, data input. \\
    6 & Data preparation of observation data to allow for comparison
with simulation results during flow model calibration.  \\
    7 & Calibration of steady state or transient flow model using existing
observations of hydraulic heads. Model calibration might be
carried out by hand (trial-and-error) or using parameter
estimation tools.  \\
    8 & Testing of the plausibility of simulated
flow velocities, flow path and water budgets.  \\
    9 & Ideally (though in
most cases not feasible) application/testing of the model for data
that have not been used during model calibration.  \\
   10 & Identification of data gaps which hamper the reliability of the
calibrated model. Preparation of plans for acquisition of
additional data, where feasible.  \\
   11 & Setup of a transport model for
advective/dispersive non-reactive transport of a single species.
Note, that the model domain of the transport model might be a
sub-domain of the flow model and that the model might have a
different dimensionality.  \\
   12 & Initial model runs of the above
single-species model to assess and fix numerical problems that
result from pure physical transport (e.g., convergence problems,
oscillations, numerical dispersion).  \\
   13 & Crude adjustment of
(longitudinal and transversal) dispersivity values to govern major
plume characteristics.  \\ \hline

  \end{tabular}
  \end{table}
  \begin{table}
  \centering
  \begin{tabular}{cp{5in}} \hline
     14 & Use of initial transport modelling
results as a (further) plausibility control for the flow model.
For example, a conservative, i.e., non-reactive species (tracer)
in the model should travel approximately at the conceptual
groundwater flow velocity. Where necessary, adjustments of the
conceptual and numerical flow model. \\
   15 & If existing reaction
modules appear unsuitable - preparation of a site-specific
reaction mod-ule, testing of the reaction module/package for
simplistic, non site-specific flow configura-tions. \\
   16 & Data preparation of observation data to allow for comparison with
simulation results during transport model calibration.  \\
   17 & If a multi-component reactive transport model is employed -- Batch-type
geochemical equi-librium modelling of background water chemistry
to assess aqueous phase - mineral equilib-ria of the
uncontaminated aquifer. Batch-type reaction simulations to assess
geochemical changes in response to the mineralisation of
hydrocarbons.  \\
   18 & Setup of multi-species, multi-component
transport model, including first estimate of reaction parameters
(from literature), where necessary.  \\
   19 & Biogeochemical transport
modelling simulations. Adjustment of reaction parameters to match
observed plume patterns - Plume length/concentrations of
hydrocarbons, redox zonation/concentrations of dissolved electron
acceptors, occurrence of reaction products.  \\
   20 & Sensitivity
analysis of calibrated model.  \\
   21 & Computation of detailed mass
budgets for dissolved species/components.  \\
   22 & Predictive model
runs to (i) estimate the long-term contaminant plume evolution and
the as-sociated risk for receptors (ii) answer `what if' - type
questions. Estimates of future stresses are needed to perform a
predictive simulation Care is needed when deciding on `future'
source conditions for predictive runs - the distri-bution and
composition of NAPL may be highly variable depending on its
weathered status and the location of the water table in the
subsurface. Note - assuming a large lateral dispersivity can lead
to predictions of short plumes due to unrealistic (over-predicted)
mixing/dilution of reactants (hydrocarbons and electron
accep-tors) and inducement of biodegradation in those zones. \\
   23 & Reporting and documenting the modelling study  \\ \hline \\
  \end{tabular}
 % \caption{Checklist of major steps in a reactive transport modelling study}
  \label{checklist}
  \end{table}